Hardening searches for dense packings by starting congruent pieces as their inscribed disks and continuously deforming them into prescribed polygons while compressing the container. Shape evolution is established in packing physics; our contribution is a controlled study of this finite-container optimization path and its failure boundaries. We compare hardening, rigid starts, growth, simulated annealing and perturbation--compression across seven families and \CsevenInstances{} instances. Gradient paths have equal step budgets; the two in-house baselines use reference-hardening pair-evaluation budgets, with measured overshoot and separate postprocessing costs. Tuned rigid starts attain the lowest container size most often overall. Hardening reaches particular square packings missed by the other tested searches within the stated run counts. A targeted higher-budget follow-up adds two such best-known successes, while identifying reused runs explicitly. At equal nominal total run budget, a mixed hardening--rigid allocation solves \PortRigidHarden{} instances versus \PortRigidSixtyFour{} for doubled rigid search; this does not demonstrate an advantage or equivalence. A preregistered pilot policy on four additional families likewise shows no detected gain. Exploratory ablations and selected replay measurements describe different search behavior without identifying a causal mechanism. Versioned plans, run outputs, saved analyses and selected replays support inspection. These results characterize an alternative search path under the tested tuning and budget design; they establish neither general competitiveness nor a validated rule for allocating a fixed compute budget.
